%%%%%%%%%%%%%%%%%%%%%%%%%%%%%%%%%%%%%%%%%%%%%%%%%%%%%
%			HLAVIČKA								%
%%%%%%%%%%%%%%%%%%%%%%%%%%%%%%%%%%%%%%%%%%%%%%%%%%%%%
\documentclass[openany,12pt]{memoir}

\usepackage[utf8]{inputenc} 
\usepackage[czech]{babel}
\usepackage[T1]{fontenc}
\usepackage[top=1.5cm, bottom=2cm, left=2cm, right=2cm]{geometry}  % --> NASTAVENÍ OKRAJŮ
\usepackage{fancyhdr}
\usepackage{graphicx}
\usepackage{lmodern}
\usepackage{xwatermark}
\usepackage{xcolor}
\usepackage{changepage}
\usepackage{newunicodechar}



%%%%%% Package na zpěvník
\usepackage[full]{leadsheets}%http://mirrors.nic.cz/tex-archive/macros/latex/contrib/leadsheets/leadsheets_en.pdf   --> dokumentace	
\definesongtitletemplate{empty}{} 
\setchords{
format = \bfseries,   %tučné akordy
minor = {mi},% 
input-notation = {german},%
output-notation = {german}%
}
\definesongtitletemplate{empty}{} 

\newlength{\drop}
\newwatermark[allpages,color=red!50,angle=0,scale=2, xpos=0,ypos=0]{\includegraphics[width=5cm]{obr/pozadi2.jpg}} %--> dvojka na pozadí


%%%% Vlastní příkazy
\newcounter{Slokočet}   %Automatické číslování slok
\newcommand{\mezera}{\vspace*{0.5cm}}   %Horizontální odsazení slok
\newcommand{\stred}{5.2cm}   %%% Na zarovnání slok doprostřed, pozn. automatičtější zarovnávání na střed nejde
\newcommand{\refren}{\mezera \noindent \textbf{R:} } %refrén
\newcommand{\sloka}{\mezera \noindent \addtocounter{Slokočet}{1} \arabic{Slokočet}. } 	%sloka, která se automaticky čísluje
\newcommand{\ssloka}{\mezera \noindent}  % vlastní číslo sloky
\newcommand{\carka}{,\:}
\newcommand{\m}[1]{\color{white}{#1}}  %Pro akordy

%%%%%%%%%%%%%%%%%%%%%%%%%%%%%%%%%%%%%%%%%%%%%%%%%%%%%%
%			NÁVOD									 %
%%%%%%%%%%%%%%%%%%%%%%%%%%%%%%%%%%%%%%%%%%%%%%%%%%%%%%
%1. Věci v hlavičce IGNOROVAT
%2. Píseň psát do prostoru mezi \begin{song} a \end{song}
%3. další řádek se značí dvěma odsazeníma (= dvakrát stisknout enter)
%4. \refren vždy na začátku refrenu a \sloka na začátek sloky (automaticky se čísluje)
%  \ = alt gr + q ; [/] = alt gr f/g ; {/} = alt gr + b/n; ^ = alt gr + 3 + mezera
%Cokoliv napíšete do ^{  } se bude brát jako akord
%když se toto bude dotýkat nějakého slova (nebude mezi tím a slovem mezera)
%tak se akordy zjeví nad slovem, ale pište to před slovo
%Když se to nedotýká slova, tak akord lítá ve vzduchu a vytiskne se větší mezera
%První možnost je asi preferovanější
%5. Akordy stačí psát jen do první sloky, když se nezmění -- kytaristi to zvládnou



%<++>
%\usepackage{subfiles}
%</++>

\begin{document}

\pagestyle{empty}
%\addtocontents{toc}{\protect\thispagestyle{empty}}
%\tableofcontents \thispagestyle{empty}\newpage
\newgeometry{top=1.5cm, bottom = 0cm, left = 2cm, right = 2cm}




%Pro kompilaci je důležitý, aby ten input a \end{document} zůstali na posledních
%dvou řádcích!!!!
\begin{song}{title=\centering Sametová \\\normalsize Žlutý pes  \vspace*{-0.3cm}}  %% sem se napíše jméno songu a autor
\moveright 4.5cm \vbox{      %Varianta č. 1  ---> Jeden sloupec zarovnaný na střed	

\sloka 
	^{G}Vzpomínám, když tehdá ^{C}před léty ^{G}začaly lítat ^{C}rakety,

	^{G}zdál se to bejt ^{C}docela dobrej ^{D}nápad,

	^{G}saxofony hrály ^{C}unyle, ^{C}frčely švédský ^{C}košile

	^{G}a někdo se moh' ^{Emi}docela dobře ^{C}flákat. ^{D}
	
\sloka
	 Když tam stál rohatej u školy, my neměli podepsaný úkoly,

	už tenkrát rozhazoval svoje sítě,

	poučen z předchozích nezdarů sestrojil elektrickou kytaru  
	
	a rock'n'roll byl zrovna narozený dítě.
	
\refren
	^{G}Vzpomínáš, takys' tu ^{D}žila, a ^{Emi}nedělej, že si ^{C}jiná,

	taková ^{G}malá pilná ^{D}včela, taková ^{C}celá ^{D}\ldots \, ^{G}sametová.
	
	
\sloka
	Přišel čas a jako náhoda byla tu bigbítová pohoda,
	
	kytičky a úsmevy sekretárok,
	
	sousedovic bejby Milena je celá blbá z Boba Dylana,
	
	ale to nevadí, já mám taky nárok.

\sloka
	Starý, mladý nebo pitomí mlátili do toho jako my,
	
	hlavu plnou Londýna nad Temží,
	
	a starej dobrej satanáš hraje u nás v hospodě mariáš
	
	a pazoura se mu trumfama jenom hemží.

\refren
	Vzpomínáš, už je to jinak, a jde z toho na mě zima,
	
	ty jsi, holka, tehdá byla taková celá \dots \,sametová.
	
\sloka
	A do toho tenhle Gorbačov, co ho znal celej Dlabačov,
	
	kopyta měl jako z Arizony,
	
	přišel a zase odešel a nikdo se kvůli tomu nevěšel
	
	a po něm tu zbyly samý volný zóny.

\refren
	Vzpomínáš, jak jsi se měla, když jsi nic nevěděla,
	
	byla to taková krásná cela a byla celá \dots
	 
	Vzpomínáš, jak jsi se měla, když jsi nic nevěděla,
	
	byla to taková krásná cela a byla celá \dots \,sametová.

}
\setcounter{Slokočet}{0}
\end{song}



\end{document}
